\chapter{Теоретические сведения}
\label{ch:theory}

\section{Система доменных имён DNS}

DNS (Domain Name System) — иерархическая распределённая система именования узлов сети \cite{rfc1034}. Пространство имён DNS имеет древовидную структуру: в корне находится безымянный корневой домен, ниже расположены домены верхнего уровня (TLD: \texttt{.ru}, \texttt{.com}, \texttt{.local} и др.), затем домены второго уровня и так далее.

Основные понятия DNS:
\begin{itemize}
    \item \textbf{Зона} --- административная единица пространства имён, за которую отвечает конкретный DNS-сервер;
    \item \textbf{Авторитативный сервер} (authoritative) --- сервер, хранящий оригинальные записи зоны и отвечающий на запросы с пометкой «авторитативный ответ»;
    \item \textbf{Рекурсивный резолвер} --- сервер, выполняющий цепочку запросов к другим DNS-серверам от имени клиента;
    \item \textbf{Форвардер} (forwarder) --- сервер, которому локальный DNS-сервер перенаправляет запросы для не локальных зон.
\end{itemize}

\section{Типы ресурсных записей DNS}

Каждая запись в зоне DNS имеет определённый тип. Основные типы, используемые в данной работе, приведены в таблице~\ref{tab:rr}.

\begin{table}[h]
    \caption{Основные типы ресурсных записей DNS}
    \label{tab:rr}
    \begin{tabular}{|l|p{10cm}|}
        \hline
        \textbf{Тип} & \textbf{Назначение} \\
        \hline
        SOA & Start of Authority --- описывает зону: первичный NS, e-mail администратора, параметры обновления \\
        \hline
        NS  & Name Server --- указывает авторитативный сервер зоны \\
        \hline
        A   & Адресная запись --- сопоставляет имя хоста с IPv4-адресом \\
        \hline
        PTR & Pointer --- обратная запись, сопоставляет IP-адрес с именем хоста \\
        \hline
        CNAME & Canonical Name --- псевдоним для другого доменного имени \\
        \hline
        MX  & Mail Exchanger --- указывает почтовый сервер домена \\
        \hline
    \end{tabular}
\end{table}

\section{Прямая и обратная зоны}

\textbf{Прямая зона} (forward zone) используется для разрешения имён в IP-адреса. Имя зоны совпадает с именем домена, например \texttt{STUDENT.GROUP.local}. Файл зоны содержит записи типов SOA, NS, A.

\textbf{Обратная зона} (reverse zone) используется для разрешения IP-адресов в имена. Имя зоны строится по принципу инвертирования первых октетов подсети с добавлением суффикса \texttt{.in-addr.arpa}. Например, для сети \texttt{192.168.N.0/24} зона будет называться \texttt{N.168.192.in-addr.arpa}. Файл зоны содержит записи типов SOA, NS, PTR.

\section{BIND9}

BIND (Berkeley Internet Name Domain) версии 9 --- наиболее распространённое и функционально полное ПО DNS-сервера \cite{bind9_admin}. В Ubuntu 20.04 устанавливается из стандартного репозитория командой \texttt{apt install bind9}. Конфигурационные файлы располагаются в каталоге \texttt{/etc/bind/}:
\begin{itemize}
    \item \texttt{named.conf} --- главный файл, включает остальные;
    \item \texttt{named.conf.options} --- глобальные настройки (форвардеры, listen-on, dnssec);
    \item \texttt{named.conf.local} --- описание локальных зон;
    \item файлы зон хранятся в \texttt{/var/lib/bind/} --- в этом каталоге BIND имеет права на запись, что необходимо для DDNS.
\end{itemize}

\section{Динамическое обновление DNS (DDNS)}

DDNS (Dynamic DNS, RFC~2136) позволяет клиентам или другим серверам автоматически обновлять записи DNS в зоне без ручного вмешательства администратора \cite{rfc2136}. В данной работе DDNS используется совместно с DHCP-сервером \texttt{isc-dhcp-server}: при получении клиентом IP-адреса DHCP-сервер автоматически создаёт A-запись и PTR-запись в DNS.

Для аутентификации обновлений используется TSIG-ключ (Transaction SIGnature) --- симметричный ключ HMAC-MD5, генерируемый утилитой \texttt{rndc-confgen}. Ключ хранится в файле \texttt{/etc/bind/rndc.key} и должен быть доступен как BIND9, так и DHCP-серверу.

\endinput
